\frfilename{mt132.tex}
\versiondate{6.4.05}
\copyrightdate{1995}

\def\chaptername{Pelengkap}
\def\sectionname{Ukuran luar dari ukuran}

\newsection{132}

Topik berikut yang ingin saya bahas adalah suatu konstruksi sederhana dengan
penerapan di mana-mana dalam teori ukuran. Dengan setiap ukuran terkait,
secara langsung, suatu ukuran luar baku (132A–132B). Jika kita mulai dengan
ukuran Lebesgue, kita kembali memperoleh ukuran luar Lebesgue (132C). Saya
memanfaatkan kesempatan ini untuk memperkenalkan gagasan `selubung terukur'
(132D–132E).

\leader{132A}{Proposisi} Misalkan $(X,\Sigma,\mu)$ suatu ruang ukur.
Definisikan $\mu^*: \Cal PX\to[0,\infty]$ dengan menuliskan

\Centerline{$\mu^*A=\inf\{\mu E:E\in\Sigma,\,A\subseteq E\}$}

\noindent untuk setiap $A\subseteq X$. Maka

(a) untuk setiap $A\subseteq X$ terdapat $E\in\Sigma$ sedemikian sehingga
$A\subseteq E$ dan $\mu^*A=\mu E$;

(b) $\mu^*$ merupakan ukuran luar pada $X$;

(c) $\mu^*E=\mu E$ untuk setiap $E\in\Sigma$;

(d) suatu himpunan $A\subseteq X$ bersifat $\mu$-terabaikan jika dan hanya
jika $\mu^*A=0$;

(e) $\mu^*(\bigcup_{n\in\Bbb N}A_n)=\lim_{n\to\infty}\mu^*A_n$ untuk
setiap barisan tak-menurun $\sequencen{A_n}$ dari subhimpunan-subhimpunan
$X$;

(f) $\mu^*A=\mu^*(A\cap F)+\mu^*(A\setminus F)$ kapan pun
$A\subseteq X$ dan $F\in\Sigma$.

\proof{{\bf (a)} Untuk setiap $n\in\Bbb N$ kita dapat memilih
$E_n\in\Sigma$ sedemikian sehingga $A\subseteq E_n$ dan
$\mu E_n\le\mu^*A+2^{-n}$; sekarang $E=\bigcap_{n\in\Bbb N}E_n\in\Sigma$,
$A\subseteq E$, dan

\Centerline{$\mu^*A\le\mu E\le\inf_{n\in\Bbb N}\mu E_n\le\mu^*A$.}

\medskip

{\bf (b)}(i) $\mu^*\emptyset=\mu\emptyset=0$. (ii) Jika
$A\subseteq B\subseteq X$, maka $\{E:A\subseteq E\in\Sigma\}\supseteq
\{E:B\subseteq E\in\Sigma\}$ sehingga $\mu^*A\le\mu^*B$.
(iii) Jika $\sequencen{A_n}$ adalah sembarang barisan dalam $\Cal PX$,
maka untuk setiap $n\in\Bbb N$ terdapat $E_n\in\Sigma$ sedemikian sehingga
$A_n\subseteq E_n$ dan $\mu E_n=\mu^*A_n$; sekarang
$\bigcup_{n\in\Bbb N}A_n\subseteq\bigcup_{n\in\Bbb N}E_n\in\Sigma$ sehingga

\Centerline{$\mu^*(\bigcup_{n\in\Bbb N}A_n)
\le\mu(\bigcup_{n\in\Bbb N}E_n)
\le\sum_{n=0}^{\infty}\mu E_n=\sum_{n=0}^{\infty}\mu^*A_n$.}

\medskip

{\bf (c)} Ini benar karena $\mu E\le\mu F$ setiap kali $E$, $F\in\Sigma$
dan $E\subseteq F$.

\medskip

{\bf (d)} Menurut (a), $\mu^*A=0$ jika dan hanya jika terdapat $E\in\Sigma$
sedemikian sehingga $A\subseteq E$ dan $\mu E=0$; tetapi inilah definisi
`himpunan terabaikan'.

\medskip

{\bf (e)} Tentu saja $\sequencen{\mu^*A_n}$ merupakan barisan tak-menurun
dengan limit paling besar $\mu^*A$, jika ditulis $A=\bigcup_{n\in\Bbb N}A_n$,
hanya karena $\mu^*B\le\mu^*C$ setiap kali $B\subseteq C\subseteq X$.
Untuk setiap $n\in\Bbb N$, ambil $E_n\in\Sigma$ sedemikian sehingga
$A_n\subseteq E_n$ dan $\mu E_n=\mu^*A_n$. Tetapkan
$F_n=\bigcap_{m\ge n}E_m$ untuk setiap $n$; maka $\sequencen{F_n}$ adalah
barisan tak-menurun dalam $\Sigma$, dan $A_n\subseteq F_n\subseteq E_n$,
sehingga $\mu^*A_n=\mu F_n$ untuk setiap $n\in\Bbb N$. Tetapkan
$F=\bigcup_{n\in\Bbb N}F_n$; maka $A\subseteq F$ sehingga

\Centerline{$\mu^*A\le\mu F=\lim_{n\to\infty}\mu
F_n=\lim_{n\to\infty}\mu^*A_n$.}

\noindent Jadi $\mu^*A=\lim_{n\to\infty}\mu^*A_n$, sebagaimana diklaim.

\medskip

{\bf (f)} Tentu saja $\mu^*A\le\mu^*(A\cap F)+\mu^*(A\setminus F)$,
menurut (b). Di sisi lain, menurut (a) terdapat $E\in\Sigma$ sedemikian
sehingga $A\subseteq E$ dan $\mu E=\mu^*A$, dan sekarang
$A\cap F\subseteq E\cap F\in\Sigma$, $A\setminus F\subseteq
E\setminus F\in\Sigma$ sehingga

\Centerline{$\mu^*(A\cap F)+\mu^*(A\setminus F)
\le\mu(E\cap F)+\mu(E\setminus F)=\mu E=\mu^*A$.}

}%end of proof of 132A

\leader{132B}{Definisi} Jika $(X,\Sigma,\mu)$ suatu ruang ukur, saya akan
menyebut $\mu^*$, sebagaimana didefinisikan dalam 132A, sebagai
{\bf ukuran luar yang diturunkan dari $\mu$}.

\cmmnt{\medskip

\noindent{\bf Catatan} Jika kita mulai dengan suatu ukuran luar $\theta$ pada
suatu himpunan $X$, membangun ukuran $\mu$ dari $\theta$ dengan metode
\Caratheodory, lalu membangun ukuran luar $\mu^*$ dari $\mu$, tidak selalu
berlaku bahwa $\mu^*=\theta$. \prooflet{\Prf\ Ambil sembarang himpunan
$X$ dengan sedikitnya tiga anggota, dan tetapkan $\theta A=0$ jika
$A=\emptyset$, $1$ jika $A=X$, dan $\bover12$ selain itu. Maka
$\dom\mu=\{\emptyset,X\}$ dan $\mu^*A=1$ untuk setiap $A\subseteq X$ tak-
kosong.\ \Qed}

Namun, masalah ini tidak muncul pada ukuran luar Lebesgue. Saya menyatakan
proposisi berikut dalam bentuk ukuran Lebesgue pada $\BbbR^r$, tetapi jika
Anda melewati \S115 saya harap Anda tetap dapat memahaminya, dan hasil-hasil
berikutnya, dalam bentuk ukuran Lebesgue pada $\Bbb R$ dengan menetapkan
$r=1$.
}%end of comment

\leader{132C}{Proposisi} Jika $\theta$ adalah ukuran luar Lebesgue pada
$\BbbR^r$ dan $\mu$ adalah ukuran Lebesgue, maka $\mu^*$, sebagaimana
didefinisikan dalam 132A, sama dengan $\theta$.

\proof{ Misalkan $A\subseteq\BbbR^r$.

\medskip

{\bf (a)} Jika $E$ terukur dan $A\subseteq E$, maka
$\theta A\le\theta E=\mu E$; jadi $\theta A\le\mu^*A$.

\medskip

{\bf (b)} Jika $\epsilon>0$, terdapat barisan
$\sequencen{I_n}$ interval setengah terbuka yang mencakup $A$, dengan
$\sum_{n=0}^{\infty}\mu I_n\le\theta A+\epsilon$ (menggunakan 114G/115G
untuk mengidentifikasi $\mu I_n$ dengan volume $\lambda I_n$ yang dipakai
dalam definisi $\theta$), sehingga

\Centerline{$\mu^*A\le\mu(\bigcup_{n\in\Bbb N}I_n)
\le\sum_{n=0}^{\infty}\mu I_n\le\theta A+\epsilon$.}

\noindent Karena $\epsilon$ sembarang, $\mu^*A\le\theta A$.
}%end of proof of 132C

\cmmnt{\medskip

\noindent{\bf Catatan} Dengan demikian, selanjutnya tidak perlu membedakan
$\theta$ dari $\mu^*$ ketika membicarakan `ukuran luar Lebesgue'. (Dalam
bahasa 132Xa di bawah, ukuran luar Lebesgue adalah `reguler'.) Secara
khusus (menggunakan 132Aa), jika $A\subseteq\BbbR^r$ terdapat himpunan
terukur $E\supseteq A$ sedemikian sehingga $\mu E=\theta A$ (bandingkan
134Fc).}%end of comment

\leader{132D}{Selubung terukur}
\cmmnt{Konsep berikut berguna dalam konteks ini.} Jika $(X,\Sigma,\mu)$
suatu ruang ukur dan $A\subseteq X$, sebuah {\bf selubung terukur} (atau
{\bf penutup terukur}) dari $A$ adalah himpunan $E\in\Sigma$ sedemikian
sehingga $A\subseteq E$ dan $\mu(F\cap E)=\mu^*(F\cap A)$ untuk setiap
$F\in\Sigma$. \cmmnt{Secara umum, tidak setiap himpunan dalam ruang ukur
memiliki selubung terukur (saya menyarankan contoh-contohnya dalam 216Yc di
Jilid 2). Namun kita memiliki hasil berikut.}

\leader{132E}{Lema} Misalkan $(X,\Sigma,\mu)$ suatu ruang ukur.

(a) Jika $A\subseteq E\in\Sigma$, maka $E$ merupakan selubung terukur dari
$A$ jika dan hanya jika $\mu F=0$ setiap kali $F\in\Sigma$ dan
$F\subseteq E\setminus A$.

(b) Jika $A\subseteq E\in\Sigma$ dan $\mu E<\infty$, maka $E$ merupakan
selubung terukur dari $A$ jika dan hanya jika $\mu E=\mu^*A$.

(c) Jika $E$ merupakan selubung terukur dari $A$ dan $H\in\Sigma$, maka
$E\cap H$ merupakan selubung terukur dari $A\cap H$.

(d) Misalkan $\sequencen{A_n}$ suatu barisan subhimpunan-subhimpunan $X$.
Andaikan setiap $A_n$ memiliki selubung terukur $E_n$. Maka
$\bigcup_{n\in\Bbb N}E_n$ merupakan selubung terukur dari
$\bigcup_{n\in\Bbb N}A_n$.

(e) Jika $A\subseteq X$ dapat dicakup oleh suatu barisan himpunan-himpunan
berukuran berhingga, maka $A$ memiliki selubung terukur.

(f) Khususnya, jika $\mu$ adalah ukuran Lebesgue pada $\BbbR^r$, maka setiap
subhimpunan $\BbbR^r$ memiliki selubung terukur untuk $\mu$.

\proof{{\bf (a)} Jika $E$ merupakan selubung terukur dari $A$, $F\in\Sigma$
dan $F\subseteq E\setminus A$, maka

\Centerline{$\mu F=\mu(F\cap E)=\mu^*(F\cap A)=0$.}

\noindent Jika $E$ bukan selubung terukur dari $A$, terdapat $H\in\Sigma$
sedemikian sehingga $\mu^*(A\cap H)<\mu(E\cap H)$. Ambil $G\in\Sigma$
sedemikian sehingga $A\cap H\subseteq G$ dan $\mu G=\mu^*(A\cap H)$,
dan tetapkan $F=E\cap H\setminus G$. Karena $\mu G<\mu(E\cap H)$, maka
$\mu F>0$; tetapi juga $F\subseteq E$ dan $F\cap A\subseteq H\cap A
\setminus G$ adalah kosong.

\medskip

{\bf (b)} Jika $E$ merupakan selubung terukur dari $A$, kita harus
memiliki

\Centerline{$\mu^*A=\mu^*(A\cap E)=\mu(E\cap E)=\mu E$.}

\noindent Jika $\mu E=\mu^*A$ dan $F\in\Sigma$ merupakan subhimpunan
$E\setminus A$, maka $A\subseteq E\setminus F$, sehingga
$\mu(E\setminus F)=\mu E$; karena $\mu E$ berhingga, diperoleh $\mu F=0$,
sehingga syarat (a) terpenuhi dan $E$ merupakan selubung terukur dari $A$.

\medskip

{\bf (c)} Jika $F\in\Sigma$ dan $F\subseteq E\cap H\setminus A$,
maka $F\subseteq E\setminus A$, sehingga menurut (a), $\mu F=0$; karena
$F$ sembarang, menurut (a) lagi, $E\cap H$ merupakan selubung terukur dari
$A\cap H$.

\medskip

{\bf (d)} Tuliskan $A$ untuk $\bigcup_{n\in\Bbb N}A_n$ dan $E$ untuk
$\bigcup_{n\in\Bbb N}E_n$. Maka $A\subseteq E$. Jika $F\in\Sigma$ dan
$F\subseteq E\setminus A$, maka, untuk setiap $n\in\Bbb N$,
$F\cap E_n\subseteq E_n\setminus A_n$, sehingga menurut (a),
$\mu(F\cap E_n)=0$. Akibatnya $F=\bigcup_{n\in\Bbb N}F\cap E_n$ adalah
terabaikan; karena $F$ sembarang, $E$ merupakan selubung terukur dari $A$.

\medskip

{\bf (e)} Ambil barisan $\sequencen{F_n}$ himpunan-himpunan berukuran
berhingga yang mencakup $A$. Untuk setiap $n\in\Bbb N$, ambil
$E_n\in\Sigma$ sedemikian sehingga $A\cap F_n\subseteq E_n$ dan
$\mu E_n=\mu^*(A\cap F_n)$ (menggunakan 132Aa); menurut (b), $E_n$ adalah
selubung terukur dari $A\cap F_n$. Menurut (d),
$\bigcup_{n\in\Bbb N}E_n$ merupakan selubung terukur dari
$\bigcup_{n\in\Bbb N}A\cap F_n=A$.

\medskip

{\bf (f)} Dalam kasus ukuran Lebesgue pada $\BbbR^r$, tentu saja
barisan $\sequencen{B_n}$ yang sama dapat digunakan untuk setiap $A$, jika
 kita mengambil $B_n$ sebagai interval setengah terbuka
$\coint{-\tbf{n},\tbf{n}}$ untuk setiap $n\in\Bbb N$, dengan menuliskan
$\tbf{n}=(n,\ldots,n)$ seperti dalam \S115.
}%end of proof of 132E

\leader{132F}{Ukuran luar penuh}\cmmnt{ Ini saat yang tepat untuk
memperkenalkan istilah berikut.} Jika $(X,\Sigma,\mu)$ suatu ruang ukur,
suatu himpunan $A\subseteq X$ disebut {\bf berukuran luar penuh} atau {\bf
tebal} jika $X$ merupakan selubung terukur dari $A$\cmmnt{; yaitu, jika
$\mu^*(F\cap A)=\mu F$ untuk setiap $F\in\Sigma$; secara ekuivalen, jika
$\mu F=0$ setiap kali $F\in\Sigma$ dan $F\subseteq X\setminus A$. Jika
$\mu X<\infty$, $A\subseteq X$ berukuran luar penuh jika dan hanya jika
$\mu^*A=\mu X$.}

\exercises{
\leader{132X}{Latihan dasar $\pmb{>}$ (a)}
%\spheader 132Xa
Misalkan $X$ suatu himpunan dan $\theta$ suatu ukuran luar pada $X$;
misalkan $\mu$ adalah ukuran pada $X$ yang didefinisikan dari $\theta$
dengan metode \Caratheodory, dan $\mu^*$ adalah ukuran luar yang diturunkan
dari $\mu$ melalui konstruksi 132A. (i) Tunjukkan bahwa $\mu^*A\ge\theta A$
untuk setiap $A\subseteq X$. (ii) $\theta$ disebut sebagai {\bf ukuran luar
reguler} jika $\theta=\mu^*$. Tunjukkan bahwa jika terdapat ukuran
$\nu$ pada $X$ sedemikian sehingga $\theta=\nu^*$, maka $\theta$ reguler.
(iii) Tunjukkan bahwa jika $\theta$ reguler dan $\sequencen{A_n}$ adalah
barisan tak-menurun dari subhimpunan-subhimpunan $X$, maka
$\theta(\bigcup_{n\in\Bbb N}A_n)=\lim_{n\to\infty}\theta A_n$.
%Rao `Carath\'eodory regular'
%132B

\spheader 132Xb Misalkan $(X,\Sigma,\mu)$ suatu ruang ukur dan $H$ sembarang
anggota $\Sigma$. Misalkan $\mu_H$ adalah ukuran subruang pada $H$ (131B)
dan $\mu^*$, $\mu_H^*$ adalah ukuran-ukuran luar yang diturunkan masing-masing
dari $\mu$, $\mu_H$. Tunjukkan bahwa
$\mu_H^*=\mu^*\restrp\Cal PH$.
%132B

\spheader 132Xc Berikan contoh ruang ukur $(X,\Sigma,\mu)$ sedemikian sehingga
ukuran $\check\mu$ yang didefinisikan dengan metode \Caratheodory dari ukuran
luar $\mu^*$ merupakan perluasan sejati dari $\mu$. \Hint{ambil $\mu X=0$.}
%132B

\sqheader 132Xd Misalkan $(X,\Sigma,\mu)$ suatu ruang ukur dan $A$ suatu
subhimpunan $X$. Andaikan $\sequencen{E_n}$ suatu barisan dalam $\Sigma$
sedemikian sehingga $\sequencen{A\cap E_n}$ saling lepas. Tunjukkan bahwa
$\mu^*(A\cap\bigcup_{n\in\Bbb N}E_n)=\sum_{n=0}^{\infty}\mu^*(A\cap E_n)$.
\Hint{gantikan $E_n$ dengan
$E_n'=E_n\setminus\bigcup_{i<n}E_i$, lalu gunakan 132Ae–132Af.}
%132B

\spheader 132Xe Misalkan $(X,\Sigma,\mu)$ suatu ruang ukur dan
$\sequencen{A_n}$ sembarang barisan subhimpunan-subhimpunan $X$. Tunjukkan
bahwa ukuran luar dari $\bigcup_{n\in\Bbb N}\bigcap_{i\ge n}A_i$ paling besar
$\liminf_{n\to\infty}\mu^*A_n$.
%132B

\spheader 132Xf Misalkan $(X,\Sigma,\mu)$ suatu ruang ukur dan anggap
$A\subseteq B\subseteq X$ sedemikian sehingga $\mu^*A=\mu^*B<\infty$.
Tunjukkan bahwa $\mu^*(A\cap E)=\mu^*(B\cap E)$ untuk setiap $E\in\Sigma$.
\Hint{selubung terukur dari $B$ merupakan selubung terukur dari $A$.}
%132E

\sqheader 132Xg Misalkan $\nu_g$ ukuran Lebesgue–Stieltjes pada $\Bbb R$,
yang dibangun seperti dalam 114Xa dari fungsi tak-menurun
$g:\Bbb R\to\Bbb R$. Tunjukkan bahwa
(i) ukuran luar $\nu_g^*$ yang diturunkan dari $\nu_g$ (132A) berimpit dengan
ukuran luar $\theta_g$ dari 114Xa;
(ii) jika $A\subseteq\Bbb R$ sembarang himpunan, maka $A$ memiliki selubung
terukur untuk ukuran $\nu_g$.
%132E

\sqheader 132Xh Misalkan $A\subseteq\BbbR^r$ suatu himpunan yang tidak
terukur menurut ukuran Lebesgue $\mu$. Tunjukkan bahwa terdapat himpunan
terukur berbatas $E$ sedemikian sehingga
$\mu^*(E\cap A)=\mu^*(E\setminus A)=\mu E>0$.
\Hint{ambil $E=E'\cap E''\cap B$, dengan $E'$ selubung terukur untuk $A$,
$E''$ selubung terukur untuk $\BbbR^r\setminus A$, dan $B$ suatu himpunan
berbatas yang sesuai.}
%132E

\spheader 132Xi Misalkan $\mu$ ukuran Lebesgue pada $\BbbR^r$ dan $\Sigma$
domainnya, serta $f$ fungsi bernilai real yang didefinisikan pada suatu
subhimpunan $\BbbR^r$ tetapi tidak $\Sigma$-terukur. Tunjukkan bahwa terdapat
$q<q'$ dalam $\Bbb Q$ dan suatu himpunan terukur berbatas $E$ sedemikian
sehingga

\Centerline{$\mu^*\{x:x\in E\cap\dom f,\,f(x)\le q\}
=\mu^*\{x:x\in E\cap\dom f,\,f(x)\ge q'\}=\mu E>0$.}

\noindent \Hint{ambil $E_q$, $E'_q$ sebagai selubung-selubung terukur untuk
$\{x:f(x)\le q\}$, $\{x:f(x)>q\}$ untuk setiap $q$. Temukan $q$ sehingga
$\mu(E_q\cap E'_q)>0$ dan $q'$ sehingga $\mu(E_q\cap E'_{q'})>0$.}
%132E

\spheader 132Xj Periksa bahwa Anda dapat mengerjakan latihan 113Yc.
%132+

\spheader 132Xk Misalkan $(X,\Sigma,\mu)$ suatu ruang ukur dan $\mu^*$ ukuran
luar yang diturunkan dari $\mu$. Tunjukkan bahwa
$\mu^*(A\cup B)+\mu^*(A\cap B)\le\mu^*A+\mu^*B$
untuk semua $A$, $B\subseteq X$.
%132B

\leader{132Y}{Latihan lanjutan (a)}
%\spheader 132Ya
Misalkan $(X,\Sigma,\mu)$ suatu ruang ukur dan $\sequencen{f_n}$ suatu
barisan fungsi bernilai real yang didefinisikan hampir di mana-mana pada
$X$. Andaikan $\sequencen{\epsilon_n}$ suatu barisan bilangan real
nonnegatif sedemikian sehingga

\Centerline{$\sum_{n=0}^{\infty}\epsilon_n<\infty$,
\quad
$\sum_{n=0}^{\infty}\mu^*\{x:|f_{n+1}(x)-f_n(x)|\ge\epsilon_n\}<\infty$.}

\noindent Tunjukkan bahwa $\lim_{n\to\infty}f_n$ terdefinisi (sebagai fungsi
bernilai real) hampir di mana-mana.
%132B

\spheader 132Yb Misalkan $(X,\Sigma,\mu)$ suatu ruang ukur, $Y$ suatu
himpunan, dan $f:X\to Y$ suatu fungsi. Misalkan $\nu$ adalah ukuran citra
$\mu f^{-1}$ (112Xf). Tunjukkan bahwa $\nu^*f[A]\ge\mu^*A$ untuk setiap
$A\subseteq X$.
%132B

\spheader 132Yc Misalkan $(X,\Sigma,\mu)$ suatu ruang ukur dengan
$\mu X<\infty$. Misalkan $\sequencen{A_n}$ suatu barisan subhimpunan-subhimpunan
$X$ sedemikian sehingga $\bigcup_{n\in\Bbb N}A_n$ berukuran luar penuh di
$X$. Tunjukkan bahwa terdapat suatu partisi $\sequencen{E_n}$ atas $X$ ke
dalam himpunan-himpunan terukur sedemikian sehingga
$\mu E_n=\mu^*(A_n\cap E_n)$ untuk setiap $n\in\Bbb N$.
%132F

\spheader 132Yd Misalkan $(X,\Sigma,\mu)$ suatu ruang ukur dan $\Cal A$ suatu
keluarga subhimpunan-subhimpunan $X$ sedemikian sehingga
$\bigcap_{n\in\Bbb N}A_n$ berukuran luar penuh untuk setiap barisan
$\sequencen{A_n}$ dalam $\Cal A$. Tunjukkan bahwa terdapat ukuran $\nu$ pada
$X$ yang memperluas $\mu$, sedemikian sehingga setiap anggota $\Cal A$
bersifat $\nu$-koterabaikan.
%132F

\spheader 132Ye Periksa bahwa Anda dapat mengerjakan latihan 113Yg–113Yh.
%132+

\spheader 132Yf\dvAnew{2010} Misalkan $(X,\Sigma,\mu)$ suatu ruang ukur.
Tunjukkan bahwa $\mu^*: \Cal PX\to[0,\infty]$ {\bf alternating
(selang-seling) untuk semua orde}, yaitu,

\Centerline{$\sum_{J\subseteq I,\#(J)\text{ genap}}
   \mu^*(A\cup\bigcup_{i\in J}A_i)
\le\sum_{J\subseteq I,\#(J)\text{ ganjil}}
\mu^*(A\cup\bigcup_{i\in J}A_i)$}

\noindent kapan pun $I$ adalah himpunan berhingga tak-kosong,
$\familyiI{A_i}$ suatu keluarga subhimpunan-subhimpunan $X$, dan $A$ suatu
subhimpunan lain dari $X$.
%mt13bits

}%end of exercises

\endnotes{
\Notesheader{132} Salah satu fakta yang hampir paling mendasar dalam teori
ukuran adalah bahwa di semua ruang ukur penting terdapat himpunan-himpunan
tak-terukur. (Untuk ukuran Lebesgue, lihat 134B di bawah.) Kita dapat
menanggapi fakta ini dengan berbagai cara. Pendekatan yang cukup berhasil
adalah mengabaikannya. Pokoknya, karena alasan yang sangat mendalam,
himpunan-himpunan dan fungsi-fungsi yang muncul dalam penerapan biasa hampir
selalu terukur, atau dapat dibuat terukur melalui manipulasi elementer; satu-
satunya pengecualian dalam matematika terapan yang saya ketahui muncul dalam
teori kendali tergeneralisasi. Sebagai matematikawan murni saya tidak nyaman
dengan pendekatan ini, dan sebagai ahli teori ukuran saya pikir pendekatan ini
menutup pintu pada beberapa gagasan paling halus dalam teori. Oleh karena itu,
dalam risalah ini himpunan-himpunan tak-terukur akan selalu hadir, meskipun
hanya secara tersirat. Dalam bagian ini saya menguraikan dua metode dasar
untuk menghadapinya: berpindah dari ukuran ke ukuran luar, yang setidaknya
memberikan semacam ukuran pada himpunan sembarang, dan gagasan `selubung
terukur', yang (ketika terdefinisi) menggambarkan wilayah tempat himpunan
tak-terukur itu harus diperhitungkan. Dalam kedua kasus kita berusaha
menguraikan himpunan tak-terukur dari luar, boleh dikatakan demikian. Tidak
ada kesulitan nyata, dan hal-hal yang perlu diperhatikan hanyalah bahwa
(i) di luar batas yang ditandai oleh 132Ee selubung terukur mungkin tidak ada
selubung terukur;
(ii) konstruksi \Caratheodory atas suatu ukuran dari ukuran luar, dan konstruksi
ukuran luar dari ukuran yang dijelaskan di sini, berkaitan erat (132C, 132Xg,
113Yc, 132Xa(i)), tetapi secara umum bukan invers satu sama lain (132B,
132Xc).
}%end of notes

\discrpage
